\documentclass{article}
\usepackage{amsmath, amssymb, cancel, mathtools, bm, xcolor, newverbs}
\usepackage[left=.5in, right=.5in, top=1in, bottom=1in]{geometry}
\usepackage[most]{tcolorbox}
\usepackage[skip=1em,indent=0pt]{parskip}
\setlength{\parindent}{0pt}

% Define custom colors
\definecolor{codebg}{HTML}{000000}    % Black background
\definecolor{codefg}{HTML}{F58727}    % Orange text

% 2. Adjust padding so it doesn't distort paragraph line spacing
\setlength{\fboxsep}{2pt} 

% 3. Clear old definition and re-route \verb properly
\let\verb\relax
\newcommand{\verb}{\collectverb{\color{codefg}\ttfamily\colorbox{codebg}}}

\begin{document}

\tableofcontents

\newpage

%%%%%%%%%%%%%%%%%%%%%%%%%%%%%%%%%%%%%%%%%%%%%%%%%%%%%%%%%%%%%%

\section{Aliases}

\subsection{Alias Commands}

\begin{itemize}
    \item \verb+alias+ - Can call at any point to see a list of available Aliases
    \item \verb+~/.zshrc+ - Location of aliases
\end{itemize}

\subsection{Helpful Aliases}

\subsubsection{Git}

\begin{itemize}
    \item \verb+alias gs='git status'+
    \item \verb+alias ga='git add .'+
    \item \verb+alias gc='git commit .'+
    \item \verb+alias gp='git push'+
    \item \verb+alias gr='git pull'+
\end{itemize}

\subsubsection{Server}

Tip: anytime I need to run something on the server, I put the code/resources in a local folder called 'server\_package', always stored in the same spot. That way I can run the aliases below to easily load those files into the server, and easily pull them back locally afterwards.

\begin{itemize}
    \item \verb+server() {ssh USERNAME@stat-$1.byu.edu}+ - type \verb+server g02+ to join the g02 server
    \item \verb+pull_server_package() {scp -r USERNAME@stat-$1.byu.edu:server_package FILE_PATH}+ 
        \begin{itemize}
            \item Type \verb+pull_server_package g02+ to retrieve the 'server\_package' folder from server g02
        \end{itemize}
    \item \verb+push_server_package() {scp -r FILE_PATH USERNAME@stat-$1.byu.edu:}+
        \begin{itemize}
            \item Type \verb+push_server_package g02+ to send the 'server\_package' folder to server g02
        \end{itemize} 
\end{itemize}

%%%%%%%%%%%%%%%%%%%%%%%%%%%%%%%%%%%%%%%%%%%%%%%%%%%%%%%%%%%%%%

\section{Servers}

\subsection{Connecting to Servers}

\verb+ssh USERNAME@stat-g01.byu.edu+ - Full line for connecting to the server (just replace the number with the corresponding server number, and USERNAME with your username)

\subsection{Moving Files}

Move file from local machine to server: \verb+scp FILE_PATH USERNAME@stat-g02.byu.edu:FILE_PATH+

Move file from server to local machine: \verb+scp USERNAME@stat-g02.byu.edu:FILE_PATH FILE_PATH+

To move a full directory, include the recursive \verb+-r+. Example: \verb+scp -r rjcass13@...+

\subsection{Running processes in background}

\subsubsection{Screens}

Running a program in a screen allows it to run in the background even when you disconnect from the server

\begin{enumerate}
    \item Login to the server
    \item Start a screen session: \verb+screen -S my_session_name+
    \item Navigate to and run the desired program
    \item Detach from the screen session: On a Mac: CTRL+A (nothing should change), then D (if it worked it will take you back to the original SSH session and say detached)
    \item To rejoin the session, while logged in to the server, do: \verb+screen -r my_session_name+. If you forgot what the session is called, do \verb+screen -ls+ to see all session names. 
    \item To end the screen session (completely kill it, no coming back), just enter \verb+exit+. If the session is frozen and the previous method doesn't work, you can use \verb+screen -d my_session_name+ in the main SSH session.
\end{enumerate}


\subsubsection{tmux}

Very similar to screen, running a program in a tmux allows it to run in the background even when you disconnect from the server. May not be included on all servers. 

\begin{enumerate}
    \item Login to the server
    \item Start a screen session: \verb+tmux new -s my_session_name+
    \item Navigate to and run the desired program
    \item Detach from the screen session: On a Mac: CTRL+B (nothing should change), then D (if it worked it will take you back to the original SSH session and say detached)
    \item To rejoin the session, while logged in to the server, do: \verb+tmux attach -t my_session+. If you forgot what the session is called, do \verb+tmux ls+ to see all tmux sessions names. 
    \item To end the screen session (completely kill it, no coming back), just enter \verb+exit+. If the session is frozen and the previous method doesn't work, you can use \verb+tmux kill-session -t session_name+ in the main SSH session.
\end{enumerate}




\end{document}